% !TeX root = ../QuestionSet.tex
% ==================== 第4章 平面向量与复数 ====================
\QuestionSetChapter{平面向量与复数}

% ===== 4.1 平面向量的概念与运算 =====
\QuestionSetSection{平面向量的概念与运算}

\questionGroup{一、选择题}

% ----------------------------------------
\begin{question}
$(1+5\text{i})\text{i}$的虚部为\choiceBox
\choices{$-1$}{0}{1}{6}
\end{question}
\begin{answer}{C}
$(1+5i)i=i-5=-5+i$，故虚部为 $1$。选 C。
\end{answer}

% ----------------------------------------
\begin{question}
帆船比赛中，运动员可借助风力计测定风速的大小和方向，测出的结果在航海学中称为视风风速，视风风速对应的向量，是真风风速对应的向量与船行风速对应的向量之和，其中船行风速对应的向量与船速对应的向量大小相等，方向相反.图1给出了部分风力等级、名称与风速大小的对应关系.已知某帆船运动员在某时刻测得的视风风速对应的向量与船速对应的向量如图2（风速的大小和向量的大小相同），单位（m/s），则真风为\choiceBox
\begin{questionImageRight}[0.28\linewidth]{image1.png}[图2]
\centering
\begin{tabular}{|c|c|c|}
\hline
等级 & 风速大小m/s & 名称 \\ \hline
2 & 1.1$\sim$3.3 & 轻风 \\ \hline
3 & 3.4$\sim$5.4 & 微风 \\ \hline
4 & 5.5$\sim$7.9 & 和风 \\ \hline
5 & 8.0$\sim$10.1 & 劲风 \\ \hline
\end{tabular}
\par\smallskip{\small 图1}
\end{questionImageRight}
\choices{轻风}{微风}{和风}{劲风}
\end{question}
\begin{answer}{A}
视风 $\vec a=(-3,-1)$，船速 $\vec b=(1,3)$，真风 $\vec n=\vec a+\vec b=(-2,2)$，$|\vec n|=\sqrt{8}=2\sqrt2\approx2.828\in[1.6,3.3)$，为轻风。选 A。
\end{answer}

% ----------------------------------------
\begin{question}[GPT生成]
已知$\vec a=(1,2)$，$\vec b=(3,-1)$，则$\vec a\cdot\vec b=$\blank{3cm}．
\end{question}
\begin{answer}{1}
$\vec a\cdot\vec b=1\cdot3+2\cdot(-1)=3-2=1$。
\end{answer}

% ===== 4.2 复数 =====
\QuestionSetSection{复数}

\questionGroup{一、选择题}

% ----------------------------------------
\begin{question}[GPT生成]
复数$z=2-3\text{i}$的共轭复数为\choiceBox
\choices{$2+3\text{i}$}{$-2-3\text{i}$}{$-2+3\text{i}$}{$3+2\text{i}$}
\end{question}
\begin{answer}{A}
复数 $a+b\text{i}$ 的共轭复数为 $a-b\text{i}$。因此 $2-3\text{i}$ 的共轭复数为 $2+3\text{i}$。选 A。
\end{answer}

% ----------------------------------------
\begin{question}[GPT生成]
若$(1+\text{i})z=2$，则$z=$\blank{3cm}．
\end{question}
\begin{answer}{$1-\text{i}$}
$z=\dfrac{2}{1+\text{i}}=\dfrac{2(1-\text{i})}{(1+\text{i})(1-\text{i})}=1-\text{i}$。
\end{answer}

\QuestionSetSection*{本章综合训练}

\questionGroup{综合解答题}

% ----------------------------------------
\begin{question}[GPT生成]
在平面直角坐标系中，已知$A(1,2)$，$B(4,0)$，$C(t,5)$。（1）求$\overrightarrow{AB}$；（2）若$\overrightarrow{AB}\perp\overrightarrow{AC}$，求$t$；（3）在（2）的条件下，求$\triangle ABC$的面积。
\end{question}
\solutionSpace{5cm}
\begin{answer}{(1) $(3,-2)$；(2) $t=3$；(3) $\dfrac{13}{2}$}
(1) $\overrightarrow{AB}=B-A=(3,-2)$，$\overrightarrow{AC}=(t-1,3)$。(2) 垂直等价于数量积为 $0$，所以 $3(t-1)+(-2)\cdot3=0$，得 $t=3$。(3) 当 $t=3$ 时，$\overrightarrow{AC}=(2,3)$，面积为 $\dfrac12|3\cdot3-(-2)\cdot2|=\dfrac{13}{2}$。
\end{answer}

% ----------------------------------------
\begin{question}[GPT生成]
设复数$z=1+\text{i}$。（1）求$z^2$；（2）求$(2-\text{i})z$；（3）求$|z^3|$，并说明能否不用直接展开$z^3$来计算。
\end{question}
\solutionSpace{5cm}
\begin{answer}{(1) $2\text{i}$；(2) $3+\text{i}$；(3) $2\sqrt2$，可以}
(1) $z^2=(1+\text{i})^2=1+2\text{i}+\text{i}^2=2\text{i}$。(2) $(2-\text{i})(1+\text{i})=2+2\text{i}-\text{i}-\text{i}^2=3+\text{i}$。(3) $|z|=\sqrt2$，故 $|z^3|=|z|^3=(\sqrt2)^3=2\sqrt2$，利用模的乘法性质即可，无需直接展开。
\end{answer}

\printChapterAnswers
